% ---
% id: reductions/rem:criteria
% kind: rem
% title: Correctness, secrecy and completeness as criteria
% summary: >-
%   Correctness and secrecy in the composable sense are functions of the output
%   state, so equality of output states transfers them. Completeness compares
%   the protocol against an honest reference channel, is not of that form and
%   is not transferred.
% prominence: standard
% status: accepted
% uses: [def:criterion]
% mentions: [def:licensed, thm:transfer, rem:completeness]
% assumes: [A-MEAS]
% refs:
%   - "portmann2022security, Eqs. (13)-(14)"
%   - "portmann2022security, Eq. (9)"
%   - "portmann2022security, Eq. (15)"
% provenance: newly-derived
% lean:
% ---

\begin{remark}\label{rem:criteria}
Correctness and secrecy in the composable sense ---
$\Pr[K_A\neq K_B]\le\varepsilon_{\rm cor}$ and
$(1-p_{\mathrm{abort}})\,D(\rho_{AE},\pi_A\otimes\rho_E)\le
\varepsilon_{\rm sec}$ of \cite[Eqs.~(13)--(14)]{portmann2022security}, in
which $\pi_A$ is the \emph{fully mixed key state} (that source writes $\tau_A$
for it; it is \emph{not} the trusted marginal $\tau_A$ of
Definition~\ref{def:licensed}, and the two never occur in one formula here)
--- are both functions of the output state in the sense of
Definition~\ref{def:criterion}, because the ideal object each compares
against is a fixed channel applied to the protocol's own output. What is
\emph{not} of that form, and is therefore not transferred by
Theorem~\ref{thm:transfer}, is a statement comparing the protocol against an
\emph{honest reference channel}: the completeness criterion of
\cite[Eq.~(9)]{portmann2022security} --- the second construction condition,
evaluated against a blank adversarial interface, from which that paper derives
robustness at its \cite[Eq.~(15)]{portmann2022security} --- and completeness
generally. Those are the concern of a different module, and
Remark~\ref{rem:completeness} records why no gap opens.
\end{remark}
